\documentclass[11pt, a4paper]{article}

% ── Packages ─────────────────────────────────────────────────────────────────
\usepackage[utf8]{inputenc}
\usepackage[T1]{fontenc}
\usepackage[spanish]{babel}
\usepackage{geometry}
\usepackage{microtype}
\usepackage{booktabs}
\usepackage{longtable}
\usepackage{array}
\usepackage{xcolor}
\usepackage{colortbl}
\usepackage{amsmath}
\usepackage{amssymb}
\usepackage{graphicx}
\usepackage{hyperref}
\usepackage{fancyhdr}
\usepackage{titlesec}
\usepackage{enumitem}
\usepackage{parskip}
\usepackage{mdframed}
\usepackage{multirow}
\usepackage{tikz}
\usetikzlibrary{shapes.geometric, arrows.meta, positioning, fit, backgrounds}

% ── Geometría ────────────────────────────────────────────────────────────────
\geometry{
  top=2.5cm, bottom=2.5cm,
  left=2.8cm, right=2.8cm,
  headheight=14pt
}

% ── Colores ──────────────────────────────────────────────────────────────────
\definecolor{tlvblue}{RGB}{31, 56, 100}
\definecolor{tlvmid}{RGB}{46, 80, 144}
\definecolor{tlvlight}{RGB}{214, 228, 248}
\definecolor{tlvgray}{RGB}{240, 244, 250}
\definecolor{tlvrule}{RGB}{200, 200, 200}
\definecolor{tlvred}{RGB}{180, 40, 40}
\definecolor{tlvgreen}{RGB}{30, 120, 60}
\definecolor{tlvorange}{RGB}{200, 100, 0}
\definecolor{structgray}{RGB}{180, 180, 180}
\definecolor{zoneyellow}{RGB}{255, 243, 200}
\definecolor{zoneblue}{RGB}{220, 235, 255}
\definecolor{zonered}{RGB}{255, 220, 220}
\definecolor{zonegreen}{RGB}{220, 245, 220}

% ── Hyperref ─────────────────────────────────────────────────────────────────
\hypersetup{
  colorlinks=true,
  linkcolor=tlvmid,
  urlcolor=tlvmid,
  pdftitle={TLV-1 DDR-002},
  pdfauthor={Programa TLV}
}

% ── Encabezados y pie ────────────────────────────────────────────────────────
\pagestyle{fancy}
\fancyhf{}
\fancyhead[L]{\small\color{tlvblue}\textbf{TLV-1} \textbar\ DDR-002}
\fancyhead[R]{\small\color{tlvblue} Edificio de Ensamblaje e Integración}
\fancyfoot[C]{\small\color{gray} Página \thepage}
\renewcommand{\headrulewidth}{0.5pt}
\renewcommand{\headrule}{\hbox to\headwidth{\color{tlvblue}\leaders\hrule height \headrulewidth\hfill}}

% ── Títulos de sección ───────────────────────────────────────────────────────
\titleformat{\section}
  {\large\bfseries\color{tlvblue}}
  {\thesection.}{0.6em}{}
  [\vspace{2pt}\color{tlvrule}\titlerule]

\titleformat{\subsection}
  {\normalsize\bfseries\color{tlvmid}}
  {\thesubsection.}{0.5em}{}

\titlespacing*{\section}{0pt}{18pt}{8pt}
\titlespacing*{\subsection}{0pt}{12pt}{4pt}

% ── Columnas de tabla ────────────────────────────────────────────────────────
\newcolumntype{L}[1]{>{\raggedright\arraybackslash}p{#1}}
\newcolumntype{C}[1]{>{\centering\arraybackslash}p{#1}}

% ── Comandos utilitarios ─────────────────────────────────────────────────────
\newcommand{\status}[1]{%
  \colorbox{tlvgreen}{\color{white}\small\bfseries\strut\ #1\ }%
}
\newcommand{\cmark}{\textcolor{tlvgreen}{\textbf{\checkmark}}}
\newcommand{\xmark}{\textcolor{tlvred}{\textbf{$\times$}}}

% ─────────────────────────────────────────────────────────────────────────────
\begin{document}

% ── PORTADA ──────────────────────────────────────────────────────────────────
\begin{titlepage}
  \centering
  \vspace*{3cm}

  {\large\color{gray}\bfseries PROGRAMA TLV}\\[0.4cm]
  {\Huge\color{tlvblue}\bfseries Test Launch Vehicle 1}\\[0.3cm]
  {\Large\color{tlvmid} DDR-002 — Design Decision Record}\\[2cm]

  \color{tlvrule}\hrule height 0.8pt
  \vspace{0.6cm}

  \begin{tabular}{L{4cm} L{9cm}}
    \rowcolor{tlvgray}
    \textbf{Documento}   & DDR-002 \\
    \textbf{Título}      & Edificio de Ensamblaje e Integración (AIT) — Configuración Horizontal \\
    \rowcolor{tlvgray}
    \textbf{Subsistema}  & Infraestructura de Tierra / Operaciones \\
    \textbf{Estado}      & \status{APROBADO} \\
    \rowcolor{tlvgray}
    \textbf{Revisión}    & 1.0 \\
    \textbf{Referencia}  & DDR-001 (Sistema de Separación de Nariz) \\
  \end{tabular}

  \vspace{0.6cm}
  \color{tlvrule}\hrule height 0.8pt

  \vfill
  {\small\color{gray} Programa TLV — Documento de Decisiones de Diseño}
\end{titlepage}

% ── 1. CONTEXTO ──────────────────────────────────────────────────────────────
\section{Contexto}

El edificio de Ensamblaje, Integración y Pruebas (AIT, por sus siglas en inglés) es la instalación donde el TLV-1 se construye, integra y verifica antes del traslado a la plataforma de lanzamiento. Dado que el TLV-1 es un vehículo de 5 m de longitud y 369 kg de masa total, la instalación puede ser considerablemente más compacta que las típicas torres de integración vertical usadas en vehículos de mayor escala.

Este DDR documenta la decisión de adoptar una configuración de ensamblaje \textbf{horizontal} y define los parámetros dimensionales, zonificación, condiciones de seguridad y secuencia de integración del edificio AIT.

% ── 2. PARÁMETROS DEL VEHÍCULO RELEVANTES ────────────────────────────────────
\section{Parámetros del Vehículo Relevantes}

\begin{longtable}{L{5.5cm} C{3cm} L{5cm}}
  \rowcolor{tlvblue}
  \color{white}\textbf{Parámetro} &
  \color{white}\textbf{Valor} &
  \color{white}\textbf{Impacto en AIT} \\
  \endfirsthead
  \rowcolor{tlvblue}
  \color{white}\textbf{Parámetro} &
  \color{white}\textbf{Valor} &
  \color{white}\textbf{Impacto en AIT} \\
  \endhead
  \rowcolor{tlvgray}
  Longitud total              & 5000 mm    & Define largo mínimo del edificio \\
  Diámetro exterior           & 330 mm     & Define ancho mínimo de la mesa \\
  Masa total al despegue      & $\approx$369 kg & Define capacidad de carga de mesa y erector \\
  \rowcolor{tlvgray}
  Masa propelente (APCP)      & $\approx$119 kg & Almacenamiento externo requerido \\
  Masa estructura estimada    & $\approx$250 kg & Carga de trabajo en mesa \\
  \rowcolor{tlvgray}
  FTS                         & Pirotécnico con pasador mecánico & Requiere zona de integración controlada \\
  Secciones principales       & 5 (nariz, control, motor, recuperación, tobera) & Define flujo de integración \\
\end{longtable}

% ── 3. PROBLEMA ──────────────────────────────────────────────────────────────
\section{Definición del Problema}

Se necesita una instalación que:

\begin{enumerate}[leftmargin=1.5cm, itemsep=3pt]
  \item Permita el ensamblaje completo del TLV-1 en condiciones controladas de temperatura, limpieza y seguridad.
  \item Minimice la complejidad constructiva y el costo de la instalación, coherente con la escala del programa.
  \item Gestione correctamente los riesgos asociados al FTS pirotécnico durante la integración.
  \item Provea espacio suficiente para el personal y el equipamiento de verificación de aviónica.
  \item Sea compatible con el traslado del vehículo a la plataforma de lanzamiento y su posterior erección vertical.
\end{enumerate}

% ── 4. ANÁLISIS DE CONFIGURACIÓN ─────────────────────────────────────────────
\section{An\'{a}lisis de Configuraci\'{o}n}

\subsection{Horizontal vs. vertical}

La integración en vertical requiere una estructura de al menos 5 m de altura interior libre, lo que implica un edificio de 7--8 m de alto con plataformas de trabajo escalonadas. Para un vehículo de esta escala, ese costo constructivo no está justificado.

La integración en horizontal es viable con las siguientes condiciones, todas resueltas:

\begin{longtable}{L{5cm} L{4.5cm} L{4cm}}
  \rowcolor{tlvblue}
  \color{white}\textbf{Preocupación} &
  \color{white}\textbf{Mitigación adoptada} &
  \color{white}\textbf{Estado} \\
  \endfirsthead
  \rowcolor{tlvblue}
  \color{white}\textbf{Preocupación} &
  \color{white}\textbf{Mitigación adoptada} &
  \color{white}\textbf{Estado} \\
  \endhead
  \rowcolor{tlvgray}
  Momento flector sobre carcasa & Mesa de soporte continuo de alta rigidez & \cmark\ Resuelto \\
  Deformación del liner por peso propio del motor & Motor se instala D$-$3 antes del lanzamiento & \cmark\ Resuelto \\
  \rowcolor{tlvgray}
  Calibración del IMU en orientación incorrecta & IMU calibrado a 0° horizontal = 90° en vertical & \cmark\ Resuelto \\
  Altura de edificio & 3 m suficiente en configuración horizontal & \cmark\ Resuelto \\
\end{longtable}

\subsection{Almacenamiento de propelente}

Los $\approx$119 kg de APCP se reciben de un proveedor externo los días previos al lanzamiento y se instalan directamente en el vehículo. El edificio AIT \textbf{no almacena propelente} en condiciones normales, lo que:

\begin{itemize}[leftmargin=1.5cm, itemsep=3pt]
  \item Elimina el requerimiento de bunker o sala de explosivos en el edificio.
  \item Simplifica radicalmente los permisos de construcción y operación.
  \item Elimina el riesgo de deterioro del propelente por almacenamiento prolongado.
  \item Reduce el personal de seguridad requerido en la instalación.
\end{itemize}

\subsection{Gestión del FTS}

El FTS pirotécnico se integra en el edificio con el pasador de seguridad mecánico instalado. El sistema permanece en estado \textbf{seguro} (pasador insertado) durante toda la permanencia en el edificio. El retiro del pasador se realiza únicamente en la plataforma de lanzamiento, en la secuencia de cuenta regresiva, por personal autorizado.

Esto clasifica el FTS como un riesgo gestionado por procedimiento, no por infraestructura, lo que no impone requerimientos constructivos adicionales al edificio.

% ── 5. DIMENSIONES Y ZONIFICACIÓN ────────────────────────────────────────────
\section{Dimensiones y Zonificación}

\subsection{Dimensiones generales}

\begin{longtable}{L{4cm} C{3cm} L{6.5cm}}
  \rowcolor{tlvblue}
  \color{white}\textbf{Dimensión} &
  \color{white}\textbf{Valor} &
  \color{white}\textbf{Justificación} \\
  \endfirsthead
  \rowcolor{tlvblue}
  \color{white}\textbf{Dimensión} &
  \color{white}\textbf{Valor} &
  \color{white}\textbf{Justificación} \\
  \endhead
  \rowcolor{tlvgray}
  Largo interior & 6.5 m & 5 m vehículo + 0.75 m espacio libre en cada extremo \\
  Ancho interior & 4.0 m & Mesa central (1.0 m) + circulación en ambos lados (1.5 m c/u) \\
  \rowcolor{tlvgray}
  Alto interior  & 3.0 m & Sin requerimiento vertical; ventilación y circulación de personal \\
  Área total     & 26 m² & Compacto, coherente con escala del programa \\
\end{longtable}

\subsection{Zonificación interior}

\begin{longtable}{C{1.2cm} L{3cm} C{2cm} L{6.5cm}}
  \rowcolor{tlvblue}
  \color{white}\textbf{Zona} &
  \color{white}\textbf{Nombre} &
  \color{white}\textbf{Área aprox.} &
  \color{white}\textbf{Función} \\
  \endfirsthead
  \rowcolor{tlvblue}
  \color{white}\textbf{Zona} &
  \color{white}\textbf{Nombre} &
  \color{white}\textbf{Área aprox.} &
  \color{white}\textbf{Función} \\
  \endhead
  \rowcolor{tlvgray}
  Z-1 & Mesa de integración & 6.5 m $\times$ 1 m & Soporte continuo del vehículo; alta rigidez estructural \\
  Z-2 & Banco de aviónica   & 2 m $\times$ 1.5 m & Verificación de AFCC, IMU, GNSS, telemetría y actuadores antes de instalación \\
  \rowcolor{tlvgray}
  Z-3 & Zona FTS            & 1.5 m $\times$ 1.5 m & Integración del FTS con pasador de seguridad; acceso restringido \\
  Z-4 & Almacén general     & 2 m $\times$ 1.5 m & Herramientas, consumibles, EPP, documentación \\
  \rowcolor{tlvgray}
  Z-5 & Circulación         & Perímetro & Paso libre mínimo de 1.2 m en ambos lados de la mesa \\
\end{longtable}

\subsection{Planta esquemática}

\begin{center}
\begin{tikzpicture}[scale=0.9]

  % Perímetro del edificio
  \draw[thick, color=tlvblue, fill=tlvgray!30]
    (0,0) rectangle (13,8);

  % Z-1: Mesa de integración (zona central longitudinal)
  \fill[zoneblue]
    (0.5, 3.2) rectangle (13, 4.8);
  \draw[tlvmid, thick]
    (0.5, 3.2) rectangle (13, 4.8);
  \node[font=\small\bfseries, color=tlvmid] at (6.75, 4.0)
    {Z-1 — Mesa de integración continua};

  % Cohete esquemático sobre la mesa
  \fill[white, draw=tlvblue, rounded corners=3pt]
    (1.0, 3.5) rectangle (12.5, 4.5);
  % Nariz
  \fill[tlvblue]
    (1.0, 3.5) -- (0.3, 4.0) -- (1.0, 4.5) -- cycle;
  % Etiqueta cohete
  \node[font=\tiny, color=tlvblue] at (6.75, 3.8) {TLV-1 (5000 mm)};

  % Z-2: Banco aviónica (esquina superior derecha)
  \fill[zonegreen]
    (9.5, 5.5) rectangle (12.5, 7.5);
  \draw[tlvgreen!70!black, thick]
    (9.5, 5.5) rectangle (12.5, 7.5);
  \node[font=\small\bfseries, color=tlvgreen!70!black, align=center] at (11.0, 6.5)
    {Z-2\\Banco\\aviónica};

  % Z-3: Zona FTS (esquina inferior derecha)
  \fill[zonered]
    (9.5, 0.5) rectangle (12.5, 2.5);
  \draw[tlvred, thick]
    (9.5, 0.5) rectangle (12.5, 2.5);
  \node[font=\small\bfseries, color=tlvred, align=center] at (11.0, 1.5)
    {Z-3\\FTS\\(acceso restringido)};

  % Z-4: Almacén (esquina superior izquierda)
  \fill[zoneyellow]
    (0.5, 5.5) rectangle (3.5, 7.5);
  \draw[tlvorange, thick]
    (0.5, 5.5) rectangle (3.5, 7.5);
  \node[font=\small\bfseries, color=tlvorange, align=center] at (2.0, 6.5)
    {Z-4\\Almacén};

  % Puerta principal (lado izquierdo inferior)
  \draw[white, line width=4pt]
    (0, 1.0) -- (0, 2.5);
  \draw[tlvblue, dashed]
    (0, 1.0) -- (0, 2.5);
  \node[font=\tiny, color=tlvblue] at (-0.8, 1.75) {Puerta};

  % Ventilacion (techo)
  \draw[tlvmid, -latex] (2, 8) -- (2, 8.6);
  \draw[tlvmid, -latex] (4.5, 8) -- (4.5, 8.6);
  \draw[tlvmid, -latex] (7, 8) -- (7, 8.6);
  \draw[tlvmid, -latex] (9.5, 8) -- (9.5, 8.6);
  \draw[tlvmid, -latex] (12, 8) -- (12, 8.6);
  \node[font=\tiny, color=tlvmid] at (6.75, 8.85) {Extraccion ventilacion forzada};

  % Cotas
  \draw[latex-latex, gray] (0, -0.4) -- (13, -0.4);
  \node[font=\small, color=gray] at (6.5, -0.75) {6.5 m (interior)};

  \draw[latex-latex, gray] (13.6, 0) -- (13.6, 8);
  \node[font=\small, color=gray, rotate=90] at (14.1, 4) {4.0 m (interior)};

\end{tikzpicture}
\end{center}

% ── 6. MESA DE INTEGRACIÓN ───────────────────────────────────────────────────
\section{Mesa de Integración (Z-1)}

La mesa es el elemento estructural crítico del edificio. Sus requerimientos son:

\begin{longtable}{L{5cm} L{8.5cm}}
  \rowcolor{tlvblue}
  \color{white}\textbf{Requerimiento} &
  \color{white}\textbf{Especificación} \\
  \endfirsthead
  \rowcolor{tlvblue}
  \color{white}\textbf{Requerimiento} &
  \color{white}\textbf{Especificación} \\
  \endhead
  \rowcolor{tlvgray}
  Longitud útil             & $\geq$ 5500 mm (vehículo + márgenes de extremo) \\
  Capacidad de carga        & $\geq$ 400 kg distribuidos (369 kg vehículo + herramientas) \\
  \rowcolor{tlvgray}
  Deflexión máxima admisible & $\leq$ 1 mm en el punto central bajo carga total \\
  Altura de trabajo         & 800--900 mm desde el suelo (ergonomía de pie) \\
  \rowcolor{tlvgray}
  Superficie                & Acero inoxidable o aluminio anodizado; no magnética en zona de aviónica \\
  Soporte del vehículo      & Perfiles en V o cunas cada 500--800 mm para distribuir la carga \\
\end{longtable}

La condición de deflexión $\leq 1\ \text{mm}$ es el constraint de diseño dominante. Para una viga simplemente apoyada de longitud $L = 5.5\ \text{m}$ con carga distribuida $q = 369/5.5 \approx 67\ \text{kg/m}$:

\begin{equation}
  \delta_{max} = \frac{5\,q\,L^4}{384\,E\,I} \leq 1\ \text{mm}
\end{equation}

El momento de inercia mínimo requerido para una mesa de acero ($E = 205\ \text{GPa}$) es:

\begin{equation}
  I_{min} = \frac{5 \times 657 \times 5.5^4}{384 \times 205 \times 10^9 \times 0.001}
  = \frac{5 \times 657 \times 915.1}{384 \times 2.05 \times 10^8}
  \approx 2.41 \times 10^{-5}\ \text{m}^4
\end{equation}

Un perfil HEB-140 ($I = 5.41 \times 10^{-5}\ \text{m}^4$) o dos perfiles UPN-120 en paralelo satisfacen este requerimiento con margen.

% ── 7. CONDICIONES AMBIENTALES Y VENTILACIÓN ─────────────────────────────────
\section{Condiciones Ambientales y Ventilación}

\begin{longtable}{L{4.5cm} L{9cm}}
  \rowcolor{tlvblue}
  \color{white}\textbf{Parámetro} &
  \color{white}\textbf{Especificación} \\
  \endfirsthead
  \rowcolor{tlvblue}
  \color{white}\textbf{Parámetro} &
  \color{white}\textbf{Especificación} \\
  \endhead
  \rowcolor{tlvgray}
  Temperatura interior      & 15--30°C (sin control activo en TLV-1; ventilación pasiva suficiente) \\
  Humedad relativa          & $\leq$ 60\% (protección de aviónica; deshumidificador portátil si necesario) \\
  \rowcolor{tlvgray}
  Ventilación               & Forzada por extractor; mínimo 6 renovaciones/hora del volumen interior \\
  Volumen interior          & $6.5 \times 4.0 \times 3.0 = 78\ \text{m}^3$ \\
  \rowcolor{tlvgray}
  Caudal mínimo de extracción & $78 \times 6 = 468\ \text{m}^3/\text{h}$ $\rightarrow$ extractor de $\geq 500\ \text{m}^3/\text{h}$ \\
  Filtración de entrada     & Prefiltro de partículas F5 (protección básica de aviónica) \\
  \rowcolor{tlvgray}
  Detectores                & CO + gases inflamables; alarma local y corte automático de ventilación \\
\end{longtable}

% ── 8. SEGURIDAD ─────────────────────────────────────────────────────────────
\section{Seguridad}

\subsection{Clasificación de riesgo del edificio}

Durante la mayor parte del ciclo de integración, el edificio contiene únicamente componentes estructurales, aviónica y el FTS asegurado mecánicamente. El propelente APCP no es almacenado en el edificio. La clasificación de riesgo es por tanto \textbf{baja} durante integración normal, elevándose a \textbf{media} únicamente durante la instalación del motor (D$-$3 a D$-$1).

\subsection{Medidas adoptadas}

\begin{longtable}{C{1cm} L{5.5cm} L{7cm}}
  \rowcolor{tlvblue}
  \color{white}\textbf{ID} &
  \color{white}\textbf{Medida} &
  \color{white}\textbf{Justificación} \\
  \endfirsthead
  \rowcolor{tlvblue}
  \color{white}\textbf{ID} &
  \color{white}\textbf{Medida} &
  \color{white}\textbf{Justificación} \\
  \endhead
  \rowcolor{tlvgray}
  S-01 & Mínimo 2 personas simultáneas durante integración de motor y FTS & Norma de seguridad para trabajo con explosivos y propelente \\
  S-02 & FTS con pasador mecánico insertado en todo momento dentro del edificio & El estado seguro es estructural, no dependiente de electrónica \\
  \rowcolor{tlvgray}
  S-03 & Acceso a Z-3 (FTS) restringido a personal autorizado & Minimiza exposición del personal no involucrado \\
  S-04 & EPP obligatorio en zona de motor: guantes, gafas, calzado antiestático & Prevención de ignición accidental por descarga electrostática \\
  \rowcolor{tlvgray}
  S-05 & Extintor de CO\textsubscript{2} (no agua ni polvo) en cada zona & El agua reacciona con APCP; el polvo contamina aviónica \\
  S-06 & Protocolo escrito de recepción del propelente (D$-$3) & Trazabilidad y verificación de integridad del grano \\
  \rowcolor{tlvgray}
  S-07 & Toma a tierra del edificio y de la mesa & Prevención de descarga electrostática durante manipulación de aviónica y FTS \\
\end{longtable}

% ── 9. SECUENCIA DE INTEGRACIÓN ──────────────────────────────────────────────
\section{Secuencia de Integración}

\begin{longtable}{C{1.2cm} L{3.5cm} C{2cm} L{6.5cm}}
  \rowcolor{tlvblue}
  \color{white}\textbf{Fase} &
  \color{white}\textbf{Actividad} &
  \color{white}\textbf{Día} &
  \color{white}\textbf{Notas} \\
  \endfirsthead
  \rowcolor{tlvblue}
  \color{white}\textbf{Fase} &
  \color{white}\textbf{Actividad} &
  \color{white}\textbf{Día} &
  \color{white}\textbf{Notas} \\
  \endhead
  \rowcolor{tlvgray}
  F-01 & Verificación aviónica (banco Z-2) & D$-$N & AFCC, IMU, GNSS, telemetría, actuadores de aletas. IMU calibrado a 0° horizontal. \\
  F-02 & Integración nariz + aviónica & D$-$N & Aviónica verificada se instala en Body\_Nariz. Pitot-estáticos conectados. \\
  \rowcolor{tlvgray}
  F-03 & Integración sección de control & D$-$N & Actuadores, servos y superficies de control. Verificación de rango de movimiento. \\
  F-04 & Integración sección de recuperación & D$-$N & Paracaídas plegado e instalado. Conexión eléctrica al AFCC. \\
  \rowcolor{tlvgray}
  F-05 & Integración FTS (Z-3) & D$-$N & Pirotécnicos del FTS instalados con pasador de seguridad. Verificación de continuidad eléctrica. \\
  F-06 & Ensayo funcional general & D$-$N & Simulación de secuencia de vuelo completa sin propelente. Telemetría end-to-end. \\
  \rowcolor{tlvgray}
  F-07 & Recepción e instalación de motor (APCP) & D$-$3 & Proveedor externo entrega el propelente. Instalación en carcasa. Verificación de igniter. \\
  F-08 & Verificación final integrada & D$-$1 & Inspección visual completa, verificación de conexiones, peso final del vehículo. \\
  \rowcolor{tlvgray}
  F-09 & Traslado a plataforma & D-0 & Vehículo completo trasladado horizontal. Erección en plataforma. Retiro de pasador FTS. \\
\end{longtable}

% ── 10. ELEMENTOS PENDIENTES ─────────────────────────────────────────────────
\section{Elementos Pendientes}

\begin{longtable}{C{1cm} L{8cm} C{2.2cm}}
  \rowcolor{tlvblue}
  \color{white}\textbf{ID} &
  \color{white}\textbf{Descripción} &
  \color{white}\textbf{Prioridad} \\
  \endfirsthead
  \rowcolor{tlvblue}
  \color{white}\textbf{ID} &
  \color{white}\textbf{Descripción} &
  \color{white}\textbf{Prioridad} \\
  \endhead
  \rowcolor{tlvgray}
  P-01 & Definir ubicación física del edificio AIT relativa a la plataforma de lanzamiento; impacta diseño del sistema de traslado. & \textcolor{tlvred}{\textbf{Alta}} \\
  P-02 & Diseñar sistema de traslado horizontal (carro o rieles) entre AIT y plataforma. & \textcolor{tlvred}{\textbf{Alta}} \\
  \rowcolor{tlvgray}
  P-03 & Diseñar erector hidráulico de pivote (DDR-003). Fuerza requerida a determinar según geometría del brazo. & \textcolor{tlvred}{\textbf{Alta}} \\
  P-04 & Verificar normativa local aplicable al almacenamiento transitorio de APCP durante F-07 y F-08 ($\approx$48 h con propelente en el edificio). & \textcolor{tlvorange}{\textbf{Media}} \\
  \rowcolor{tlvgray}
  P-05 & Especificar mesa de integración: selección de perfil estructural, material de superficie y sistema de cunas de soporte del vehículo. & \textcolor{tlvorange}{\textbf{Media}} \\
  P-06 & Definir sistema de toma a tierra y protección antiestática de la mesa. & \textcolor{tlvorange}{\textbf{Media}} \\
  \rowcolor{tlvgray}
  P-07 & Elaborar protocolo escrito de recepción de propelente (S-06). & \textbf{Baja} \\
\end{longtable}

\end{document}